% page 263 du fac-simile (image p-262 du scan)
% bloc 170.2 x 237.7 mm ; marge gauche 31.8 mm
\pgc{10.160mm}{0.000mm}{
\l{\cel{52}255}
\l{}
\l{\sou{kalko}. (\sou{kemio}).Kali-oxido, alkalio minerala. - DEFIRS.}
\l{}
\l{\sou{kalkane}o.(\sou{anat.}) Osto di la tarso, qua formacas la ta-}
\l{\cel{3}lono. - EF.}
\l{}
\l{\sou{kalkolo}. (\sou{patol.}) Konkreciono stonatra qua naskas acidente en}
\l{\cel{3}ula organi (hepato, veziko, e c.) - EFIS.}
\l{}
\l{\sou{kalkular}. (\sou{trans.)}\cel{3}I. Determinar, per operaci matematikala,}
\l{\cel{3}ye nombri donita, nombron quan onu serchas. - II. (\sou{metaf.})}
\l{\cel{3}Agar sua aranji, kombinar la kozi, egarde di la skopo atin-}
\l{\cel{3}genda. - DEFIRS.}
\l{}
\l{\sou{kalma}. Stando di to quo indijas agiteso. - EFIS.}
\l{}
\l{\sou{kalmar}o.(\sou{zool.)}\cel{2}Molusko ek la genero \textquotedbl{}sepii\textquotedbl{}, qua sekrecas}
\l{\cel{3}liquido nigra quan lu disjetas cirkum su kande onu atakas lu.}
\l{\cel{3}- L.\cel{1}\sou{loligo.}\cel{1}- DEFIRS.}
\l{}
\l{\sou{kalomelo}. (\sou{farmacio}).Sub-kloruro di merkurio, korpo blanka, ne-}
\l{\cel{3}dissolvebla en aquo, uzata medicine kom purgivo, e kom}
\l{\cel{3}kontre-vermo. - DEFIRS.}
\l{}
\l{\sou{kaloro}. Kauzo fizikala di ula modifiki molekulala di la korpo,}
\l{\cel{3}vario da volumino - dilateso ; vario di stando - liquidesko,}
\l{\cel{3}vaporesko). - EFIS.}
\l{}
\l{kalorio.\cel{2}(\sou{fiziko)}. Mezur-unajo pri kaloro : la quanto de kaloro}
\l{\cel{3}qua esas necesa por plualtigar ye un grado centigrada la tem-}
\l{\cel{3}peraturo di un kilogramo de aquo. - DEFIRS.}
\l{}
\l{kaloriko.\cel{3}Kaloro-principo quan kontenas la korpi. - EFIS.}
\l{}
\l{\sou{kalorimetro}. Fizik-instrumento, uzata por determinar la kaloro}
\l{\cel{3}specifika di la korpi. - DEFIS.}
\l{}
\l{kalorimetrio.\cel{2}La cienco qua traktas la mezuro di la kalorii. -}
\l{\cel{3}DEFIS.}
\l{}
\l{kaloto. I. Mikra boneto ek drapo, veluro, e c., ronda, qua kovras}
\l{\cel{3}la suprajo di la kapo. - II. (\sou{geom.}) En sfero quan sekas}
\l{\cel{3}plano : la surfaco di un ek la du parti, precipue di la plu}
\l{\cel{3}mikra. - III. Quaza kapuco ek stano, ye qua onu envelopas la}
\l{\cel{3}stopilo e la kolo di la flakoni,di la liquor-boteli, e c. - DEFI}
\l{}
\l{kalquar.\cel{2}(trans.) Kopiar desegnuro tra papero-folio diafana,}
\l{\cel{3}per streki kontinua. - DFIS.}
\l{}
\l{kalsono.\cel{2}Suba vesto, quaza bracho, ek telo, flanelo, e c. - FRS.}
\l{}
\l{kalumniar.\cel{2}(trans.)Dicar neexaktaji mala pri (ulu). - EFIS.}
\l{}
\l{kalva. Qua havas nur tre poka hari. - DEFIS.}
}
